\lyseis{Β3.4}{Είδη τριγώνων και στοιχεία τους}
\sol{1} α) Σωστό. β) Λάθος: για να είναι οξυγώνιο πρέπει και οι τρεις γωνίες να είναι οξείες. γ) Λάθος: δύο πλευρές ίσες, όχι όλες. δ) Σωστό: όλες οι πλευρές είναι ίσες, οποιαδήποτε μπορεί να είναι βάση. ε) Λάθος: έχει το πολύ μία αμβλεία γωνία.
\sol{2} α) ισόπλευρο \quad β) ισοσκελές \quad γ) σκαληνό.
\sol{3} α) οξυγώνιο \quad β) ορθογώνιο \quad γ) αμβλυγώνιο.
\sol{4} Ελεύθερη κατασκευή. Στο οξυγώνιο και τα τρία ύψη είναι εντός του τριγώνου. Στο ορθογώνιο οι δύο κάθετες πλευρές είναι ύψη. Στο αμβλυγώνιο δύο ύψη βρίσκονται εκτός του τριγώνου.
\sol{5} Ελεύθερη κατασκευή. Οι τρεις διάμεσοι περνούν από το ίδιο σημείο, που βρίσκεται μέσα στο τρίγωνο.
\sol{6} Ελεύθερη κατασκευή. Οι δύο γωνίες είναι ίσες, γιατί η διχοτόμος χωρίζει τη γωνία στα δύο.
\sol{7} Το $\mathrm{Μ}$ είναι το μέσο της $\mathrm{ΒΓ}$, άρα $\mathrm{ΒΓ}=2\cdot2{,}6=5{,}2$ cm.
\sol{8} Ισοσκελές (δύο ίσες πλευρές), με βάση $40$ cm. Περίμετρος $30+30+40=100$ cm.
\sol{9} Περίμετρος $3\cdot60=180$ cm.
\sol{10} Η πλευρά είναι $24:3=8$ μ.
\sol{11} Οι δύο ίσες πλευρές έχουν μαζί $50-20=30$ cm, άρα η καθεμία είναι $30:2=15$ cm.
\sol{12} Βάση είναι η πλευρά που δεν είναι ίση με τις άλλες, δηλαδή τα $8$ μ. Περίμετρος $5{,}2+5{,}2+8=18{,}4$ μ.
\sol{13} Υπάρχουν: σκαληνό οξυγώνιο, σκαληνό ορθογώνιο, σκαληνό αμβλυγώνιο, ισοσκελές οξυγώνιο, ισοσκελές ορθογώνιο, ισοσκελές αμβλυγώνιο και ισόπλευρο οξυγώνιο. Ισόπλευρο ορθογώνιο ή αμβλυγώνιο δεν υπάρχει, γιατί όλες οι γωνίες του ισόπλευρου είναι ίσες (και $60^\circ$), άρα όλες οξείες. Ελεύθερα σχέδια για τα παραδείγματα.
\sol{14} Σε οξυγώνιο, σε ορθογώνιο ή σε αμβλυγώνιο (π.χ. ισοσκελές με γωνία βάσης $45^\circ$ είναι ορθογώνιο).
\sol{15} Όχι. Ένα ισόπλευρο τρίγωνο έχει όλες τις γωνίες ίσες (και ίσες με $60^\circ$), άρα δεν μπορεί να έχει ορθή γωνία.
\sol{16} Δύο, στο αμβλυγώνιο τρίγωνο. Στο οξυγώνιο κανένα, στο ορθογώνιο κανένα (οι κάθετες πλευρές είναι ύψη πάνω στις πλευρές).
\sol{17} Ελεύθερη απάντηση. Π.χ. φέτα πίτσας: ισοσκελές, πινακίδα «Δώσε προτεραιότητα»: ισόπλευρο.
\sol{18} Ελεύθερη κατασκευή. Διπλώνουμε ώστε οι δύο άκρες της πλευράς να συμπέσουν: η πτυχή τέμνει την πλευρά στο μέσο της. Η διάμεσος ενώνει το μέσο με την απέναντι κορυφή.
